\documentclass[11pt,a4paper]{article}
\usepackage{fontspec,amsmath,amssymb,amsthm,microtype}
\setmainfont{lmroman10-regular.otf}[BoldFont=lmroman10-bold.otf,ItalicFont=lmroman10-italic.otf,BoldItalicFont=lmroman10-bolditalic.otf]
\usepackage{xcolor}
\usepackage[most]{tcolorbox}
\usepackage{mathrsfs}
\providecommand{\lt}{<}
\providecommand{\gt}{>}
\setmonofont{DejaVuSansMono.ttf}[BoldFont=DejaVuSansMono-Bold.ttf,ItalicFont=DejaVuSansMono-Oblique.ttf,BoldItalicFont=DejaVuSansMono-BoldOblique.ttf,Scale=0.86]
\newfontfamily\leanfallback{STIXTwoMath-Regular.otf}[Scale=0.86]
\definecolor{leanrule}{gray}{0.62}
\definecolor{leanback}{gray}{0.962}
\definecolor{leantitle}{gray}{0.22}
\definecolor{leanmuted}{gray}{0.38}
\newtcolorbox{leanbox}[1]{enhanced,breakable,colback=leanback,colframe=leanrule,
  boxrule=0.4pt,arc=0.6mm,left=2.4mm,right=2.4mm,top=1.8mm,bottom=1.8mm,
  before skip=0.8em,after skip=1em,fonttitle=\small\bfseries,coltitle=leantitle,
  colbacktitle=leanback,titlerule=0.3pt,title={#1}}
\newlength{\leancharwidth}
\newenvironment{leancode}{\par\smallskip\begingroup\ttfamily\footnotesize\raggedright
  \settowidth{\leancharwidth}{M}\setlength{\parskip}{0pt}\setlength{\parindent}{0pt}
  \frenchspacing}{\par\endgroup\smallskip}
\newcommand{\leanline}[2]{\hangindent=\dimexpr#1\leancharwidth+4\leancharwidth\relax
  \hangafter=1\noindent\hspace*{#1\leancharwidth}#2\par}
\newcommand{\leansource}[1]{\par{\footnotesize\color{leanmuted}#1}\par}
\newcommand{\leanlabel}[1]{\par\smallskip\noindent{\small\bfseries #1}\par}
\usepackage[margin=25mm]{geometry}
\usepackage[unicode,hidelinks]{hyperref}
\hypersetup{pdfauthor={LangevinLean Project}}
\setlength{\emergencystretch}{3em}
\setlength{\parskip}{0.4em}
\setlength{\parindent}{0pt}
\allowdisplaybreaks[2]
\raggedbottom
\title{\texorpdfstring{A finite coefficient for the splitting error floor}{A finite coefficient for the splitting error floor}}
\author{LangevinLean Project}
\date{}
\begin{document}\maketitle
Draft mathematical exposition; these notes may contain mistakes.


The explicit local forcing is at most \(mh^2C\) when its potential gap is at most \(hK^2/2\). Its coefficient uses curvature, smooth-gradient regularity, drift and noise moments, and the bounded perturbation-gradient magnitude. It contains no inverse envelope gradient constant. Thus the one-step squared forcing is \(O(h^2)\) under \(\delta\le hK^2/2\), as required for a transport error tending to zero with the step.


\section{\texorpdfstring{The explicit nonnegative coefficient}{The explicit nonnegative coefficient}}


For \(m>0\) and \(G,K\ge0\) define

\[C(m,L,G,\sigma,K)=\frac{4L^2}{m^2}\left(\frac{G^2}{m}+\sigma^2\right)+\frac2m(GK+K^2).\]

Every squared term is nonnegative and the denominators are positive, giving \(C\ge0\). Each summand of \(C\) is finite under \(m>0\) and the stated nonnegative magnitude bounds. The coefficient is not asserted optimal; its explicit dependence on dimension, curvature, and particle factors is retained in the quantitative estimate.

\begin{leanbox}{Lean: \texttt{splittingRMSConstant}, \texttt{splittingRMSConstant\_\allowbreak{}nonneg}}
\leansource{\texttt{splittingRMSConstant} \quad ProximalLangevin/\allowbreak{}SplittingRate.lean, lines 17–18}
\begin{leancode}
\leanline{0}{def\ splittingRMSConstant\ (m\ L\ G\ noise\ K\ :\ ℝ)\ :\ ℝ\ :=}
\leanline{2}{4\ *\ L\ \textasciicircum{}\ 2\ /\ m\ \textasciicircum{}\ 2\ *\ (G\ \textasciicircum{}\ 2\ /\ m\ +\ noise\ \textasciicircum{}\ 2)\ +\ 2\ /\ m\ *\ (G\ *\ K\ +\ K\ \textasciicircum{}\ 2)}
\end{leancode}
\leansource{\texttt{splittingRMSConstant\_\allowbreak{}nonneg} \quad ProximalLangevin/\allowbreak{}SplittingRate.lean, lines 21–25}
\begin{leancode}
\leanline{0}{theorem\ splittingRMSConstant\_nonneg\ \{m\ L\ G\ noise\ K\ :\ ℝ\}}
\leanline{4}{(hm\ :\ 0\ <\ m)\ (hG\ :\ 0\ ≤\ G)\ (hK\ :\ 0\ ≤\ K)\ :}
\leanline{4}{0\ ≤\ splittingRMSConstant\ m\ L\ G\ noise\ K}
\end{leancode}
\leanlabel{How the Lean proof follows the argument}
Formal correspondence: splittingRMSConstant : A finite coefficient for the squared splitting error floor. splittingRMSConstant\_nonneg : The coefficient is nonnegative under the actual moment-bound hypotheses.



\texttt{splitting\allowbreak{}RMSConstant} is a definition with defining expression \texttt{4 * L \textasciicircum{} 2 /\allowbreak{} m \textasciicircum{} 2 * (G \textasciicircum{} 2 /\allowbreak{} m + noise \textasciicircum{} 2) + 2 /\allowbreak{} m * (G * K + K \textasciicircum{} 2)}. The right-hand side specifies the object used in the mathematical statement.



\texttt{splitting\allowbreak{}RMSConstant\_\allowbreak{}nonneg}: the proof is the application \texttt{unfold splitting\allowbreak{}RMSConstant}. It specializes that result with the hypotheses and objects in the exact statement.
\end{leanbox}

\section{\texorpdfstring{The quadratic-step bound on local forcing}{The quadratic-step bound on local forcing}}


Assume \(h\le1/m\) and the stated gap bound. The estimates \((Gh+\sigma\sqrt h)^2\le2G^2h^2+2\sigma^2h\le2h(G^2/m+\sigma^2)\) and \(h+1/m\le2/m\) bound the weighted smooth defect by \(4L^2h^2(G^2/m+\sigma^2)/m\). The gap contributes \(2h\delta\le h^2K^2\). Adding \(2h^2GK+h^2K^2\) proves \(\mathcal F\le mh^2C\), whose geometric sum has floor \(hC\).



For the expression \((Gh+\sigma\sqrt h)^2\), apply \((u+v)^2\le2u^2+2v^2\) and \(h\le1/m\) to obtain

\[(Gh+\sigma\sqrt h)^2\le2h(G^2/m+\sigma^2).\]

The remaining weighted factor satisfies \((1+(mh)^{-1})h^2\cdot2h=2h^2(h+1/m)\le4h^2/m\). The gap bound is \(2h\delta\le h^2K^2\); addition of \(h^2K^2\) and \(2h^2GK\) gives exactly \(mh^2\,C\). Positivity \(m,h>0\) justifies all denominators.



The use of \((a+b)^2\le2a^2+2b^2\) separates drift and noise. Under \(h\le1/m\) the additional powers of \(h\) are bounded by curvature constants, leaving a coefficient times \(h^2\). The gap contributes \(h^2K^2\), and the sum is \(mh^2C\). Division by the later contraction gap \(mh\) gives a squared transport error linear in the step.

\begin{leanbox}{Lean: \texttt{splittingCouplingForcing\_\allowbreak{}le}}
\leansource{\texttt{splittingCouplingForcing\_\allowbreak{}le} \quad ProximalLangevin/\allowbreak{}SplittingRate.lean, lines 29–78}
\begin{leancode}
\leanline{0}{theorem\ splittingCouplingForcing\_le\ \{m\ h\ L\ G\ noise\ K\ delta\ :\ ℝ\}}
\leanline{4}{(hm\ :\ 0\ <\ m)\ (hh\ :\ 0\ <\ h)\ (hsmall\ :\ h\ ≤\ 1\ /\ m)}
\leanline{4}{(\_hG\ :\ 0\ ≤\ G)\ (\_hnoise\ :\ 0\ ≤\ noise)\ (\_hK\ :\ 0\ ≤\ K)}
\leanline{4}{(hdelta\ :\ delta\ ≤\ h\ *\ K\ \textasciicircum{}\ 2\ /\ 2)\ :}
\leanline{4}{splittingCouplingForcing\ m\ h\ L\ G\ noise\ K\ delta\ ≤}
\leanline{6}{m\ *\ h\ \textasciicircum{}\ 2\ *\ splittingRMSConstant\ m\ L\ G\ noise\ K}
\end{leancode}
\leanlabel{How the Lean proof follows the argument}
\texttt{splitting\allowbreak{}Coupling\allowbreak{}Forcing\_\allowbreak{}le} is the scalar estimate used after the stationary coupling has supplied its four forcing terms. \texttt{hlocal} bounds the squared increment by twice the sum of squares. \texttt{h\allowbreak{}Gscale} uses \texttt{hsmall} to remove one time power, and \texttt{hlocal'} collects that estimate. \texttt{hweight} allows multiplication by the Young weight. The identity \texttt{heq} and bound \texttt{htwice} control its coefficient; \texttt{hdelta'} inserts the potential gap. Finally \texttt{h\allowbreak{}C} expands the proposed constant and \texttt{hfirst} bounds the weighted smooth-error term. Unfolding \texttt{splitting\allowbreak{}Coupling\allowbreak{}Forcing} leaves precisely the scalar inequality stated above. The actual law recurrence is established separately in the splitting endpoint module.



Formal correspondence: splittingCouplingForcing\_le : With a gap of at most \(hK^2/2\), the derived one-step forcing is bounded by \(mh^2C\), where \(C\) is independent of the smoothing parameter.
\end{leanbox}

\end{document}
